\subsection{A formal leading-coefficient criterion}\label{sec:stationaryballs}

For a polynomial phase, the remaining shrinking argument can be
replaced by a valuation calculation. The construction below includes
critical points that collide, remain multiple, or escape to infinity.
The key is to group them before making a Gaussian approximation.

Write $v$ for the valuation of a formal parameter $r$, normalized by
$v(r)=1$. The valuation of a polynomial is the least valuation of its
coefficients. We use the formal Newton--Puiseux theorem: finitely many
algebraic elements over $\C((r))$ belong to one field
$\C((r^{1/N}))$. Convergence is not required~\cite[Chapter~2]{Wall}.

\begin{lemma}[Polynomial separation of formal balls]\label{lem:ballseparation}
Let finitely many disjoint balls have centers $c_B$ in a Puiseux field
and radii specified by $v(x-c_B)\geq\rho_B$. After one common
ramification, put $\delta_B=r^{\rho_B}$. For any one ball $B$ and
integer $j\geq0$, there is a polynomial $q_{B,j}(x)$ over the Laurent
field whose substitution $x=c_C+\delta_C y$ has integral coefficients
for every ball $C$, and whose reduction is $y^j$ on $B$ and zero on
every other ball.
\end{lemma}
\begin{proof}
The finitely many valuations $v(c_B-c_C)$ give a finite tree of
clusters of centers. At a node with scale $r^d$, the coordinate
$u=(x-c)/r^d$ has distinct constant reductions $a_1,\ldots,a_s$
on its children. The Lagrange polynomial
\[
 e_i(u)=\prod_{k\ne i}\frac{u-a_k}{a_i-a_k}
\]
reduces to one on child $i$ and zero on every other child.

Start at the chosen leaf with $((x-c_B)/\delta_B)^j$. Induct upward.
Suppose a polynomial $h$ already has the required reductions on the
chosen child. On each of the other child balls its coefficients have
some finite lower valuation. The restriction of $e_i$ there has
strictly positive valuation. Consequently $e_i^M h$, for a sufficiently
large fixed integer $M$, has positive coefficient valuation on all
those children. On the chosen child $e_i^M=1+$ positive-order terms,
so the existing reductions are unchanged. There are finitely many
nodes, and this constructs the required polynomial. A node with just
one leaf requires no filter. This also explains why independently
normalized Lagrange polynomials, with no control of their growth on
other scales, would not suffice.
\end{proof}

\begin{theorem}[Formal stationary balls]\label{thm:stationaryballs}
Let $\varphi(r,x)\in\C((r))[x]$ have degree $D\geq2$, and let
$h(r)\in\C((r))\setminus\{0\}$. Put
\[
 M=\C((r))[x]/(\partial_x+\varphi_x)\C((r))[x],\qquad
 \nabla[q]=[q_r+(\varphi_r+h'/h)q].
\]
After a finite ramification, this module has a full frame of rank-one
formal exponential functionals. If $a_D$ is the leading coefficient
of $\varphi$, every zero-exponential functional in this frame takes
the seed $[1]$ to a series of valuation at least
\begin{equation}\label{eq:valuationbound}
 v(h)-\frac{v(a_D)}D.
\end{equation}
In particular, if $Dv(h)>v(a_D)$, the differential submodule generated
by $[1]$ has no horizontal quotient taking $[1]$ to the constant one.
\end{theorem}
\begin{proof}
The two twisted derivatives commute, so the parameter connection
is well defined on the quotient. A nonzero polynomial of degree $m$
has spatial twisted derivative of degree $m+D-1$. Reduction of
successive leading terms therefore gives a basis
$1,x,\ldots,x^{D-2}$ and dimension $D-1$.

Split $\varphi_x$, retaining root multiplicities. For a critical root
$\xi$, expand
\[
 \varphi(\xi+u)=c_\xi+\sum_{j=2}^D a_{\xi,j}u^j,
 \qquad
 \rho_\xi=\max_{2\leq j\leq D,\ a_{\xi,j}\ne0}
                    \frac{-v(a_{\xi,j})}{j}.
\]
Enlarge the ramification so that all these scales are integral, and
set $\delta_\xi=r^{\rho_\xi}$. The polynomial
\[
 \varphi(\xi+\delta_\xi y)-c_\xi
\]
has integral coefficients, no linear term, and at least one unit
coefficient. Its reduction $F_\xi(y)$ has degree $d_\xi\geq2$.
Here and below we may replace $r$ by a root of $r$ and renormalize
the valuation; every displayed strict inequality is unchanged.
As $a_{\xi,D}=a_D$,
\begin{equation}\label{eq:balllower}
 \rho_\xi\geq-v(a_D)/D.
\end{equation}

\emph{The balls partition all critical roots.}
If another critical root $\zeta$ satisfies
$v(\zeta-\xi)\geq\rho_\xi$, translate the integral scaled polynomial
by $(\zeta-\xi)/\delta_\xi$. All its coefficients remain integral,
and the coefficient of degree $d_\xi$ remains a unit: higher-degree
terms have positive valuation. The linear term vanishes because
$\zeta$ is critical. It follows from the definition of the maximum
that $\rho_\zeta=\rho_\xi$. Thus their balls are identical.
If two of the balls intersect, one contains the other; the same
argument makes them identical. Otherwise they are disjoint.

The scaled derivative has reduction $F_\xi'$ of degree $d_\xi-1$.
It therefore has exactly $d_\xi-1$ integral roots, with multiplicity.
Here is a direct check of that count: factor it in a Puiseux field.
An integral root contributes a linear factor to the reduction;
for a root of negative valuation, factor out that root and the
remaining factor reduces to a constant. Multiplicativity of the
coefficient valuation leaves a nonzero constant scalar. Hence the
degree of the reduction counts exactly the integral roots.
The ball $B$ thus contains $d_B-1$ critical roots. Since
$\deg\varphi_x=D-1$,
\begin{equation}\label{eq:ballrankcount}
 \sum_B(d_B-1)=D-1.
\end{equation}

\emph{A functional for every dimension in a ball.}
The fixed polynomial $F_B$ has twisted quotient
\[
 H_B=\C[y]/(\partial_y+F_B')\C[y]
\]
with basis $1,y,\ldots,y^{d_B-2}$. Let $\ell^0_{B,j}$ be its dual
basis. Write
\[
 \varphi(\xi_B+\delta_B y)=c_B+F_B(y)+E_B(r,y),
 \qquad E_B\in r\C[[r]][y].
\]
For a polynomial amplitude $q$ define
\begin{equation}\label{eq:ballfunctional}
 \ell_{B,j}(q)=h\delta_B e^{c_B}
   \ell^0_{B,j}\!\left(e^{E_B(r,y)}q(\xi_B+\delta_B y)\right).
\end{equation}
Each coefficient uses finitely many polynomial terms. The exponential
in this formula is purely formal. The product rule conjugates the
spatial twisted derivative to $\partial_y+F_B'$, so the functional
annihilates the defining relations of $M$.

We spell out the parameter check because it retains the moving volume.
At fixed $y$, put $w=\xi_B'+\delta_B'y$. Differentiating
\eqref{eq:ballfunctional}, then subtracting its value on
$q_r+(\varphi_r+h'/h)q$, leaves, inside the fixed functional,
\[
 e^{E_B}\left(wq_x+\frac{\delta_B'}{\delta_B}q+\varphi_xwq\right)
 =\frac1{\delta_B}(\partial_y+F_B')
                      \left(e^{E_B}wq\right).
\]
It is zero. Thus the functionals are horizontal, including the
Jacobian $\delta_B$ and normalization $h$.

\emph{Completeness, not just a rank count.}
Choose the polynomials $q_{B,j}$ of
Lemma~\ref{lem:ballseparation}, for $0\leq j<d_B-1$.
Divide each evaluation column by its scalar $h\delta_Be^{c_B}$.
The resulting matrix reduces to the identity: $E_B$ has positive
order and the reductions of the amplitudes are the chosen monomials
on one ball and zero elsewhere. The matrix is therefore invertible.
Leading-degree reduction gives $\dim M=D-1$, so
\eqref{eq:ballrankcount} accounts for the entire module.

Remove from $c_B$ its finite negative-power part $q_B(r)$.
The remaining exponential is an invertible formal power series,
up to a nonzero constant. Formula~\eqref{eq:ballfunctional} and the
invertible matrix give the asserted direct sum of rank-one factors
$e^{q_B}$ times powers of $r$. If $q_B=0$, its value on the seed
has valuation at least $v(h)+\rho_B$, proving
\eqref{eq:valuationbound}.

Finally, after ramification this is a semisimple differential module.
A quotient on the submodule generated by the seed extends to the
whole module. Nonzero exponential factors have no horizontal map
to the Laurent field, because the logarithmic derivative of a
Laurent series cannot be the derivative of a nonzero negative-power
exponent. On the zero part, every horizontal functional is a
constant combination of those in the full frame. If the bound is
strictly positive, their seed values cannot sum to one.
\end{proof}

The construction is a formal stationary-phase argument. Its classical
background includes Newton--Puiseux and the formal stationary-phase
theory of exponential connections; Sabbah's local Fourier--Laplace
formula is a primary reference for the latter~\cite{SabbahStationary}.
We do not identify that formula's scope with the particular proof
above or claim priority for this formulation. The proof here also
allows nonconvergent formal Laurent coefficients. It uses a valued
field, not an arbitrary coefficient ring.

\begin{corollary}[Leading-coefficient obstruction]\label{cor:leadingobstruction}
For a rational polynomial phase $\varphi(t,x)$ of degree $D\geq2$
and nonzero rational scalar normalization $h(t)$, let $v_\infty$
mean order at $\varepsilon=1/t=0$. If
\begin{equation}\label{eq:leadingcriterion}
 v_\infty(h^D/a_D)>0,
\end{equation}
its seed submodule cannot have a rational horizontal quotient with
seed value one. In particular any horizontal formal functional at
another point cannot take that seed to one.
\end{corollary}
\begin{proof}
Apply the theorem at infinity. A rational horizontal quotient remains
a quotient on extension to the formal Laurent field. A formal
functional with seed value one restricts to such a rational quotient:
its values on successive seed derivatives are $1,0,0,\ldots$.
\end{proof}

The inequality must be strict. For example
$\varphi=-x^2/(2t^2)$ and $h=1/t$ have
$v_\infty(h^2/a_2)=0$, and their normalized Gaussian period is a
nonzero constant. Thus the earlier adverse scaling example is also
an exact boundary test of the criterion.

